\subsubsection{Simon as a Fourier / Period-Finding Algorithm}\label{sec:simon-as-a-fourier-period-finding-algorithm}

We've been describing Simon using Hadamard gates $H^{\otimes n}$, but there is a deeper interpretation: \hl{$H^{\otimes n}$ is the quantum Fourier transform appropriate to $\mathbb{Z}_2^n$}.

\highspace
\begin{flushleft}
    \textcolor{Green3}{\faIcon{question-circle} \textbf{What is $\mathbb{Z}_2^n$?}}
\end{flushleft}
It is essentially the mathematical structure of our $n$-bit strings:
\begin{equation*}
    \left\{0, 1\right\}^{n}
\end{equation*}
For example, for $n=2$, we have:
\begin{equation*}
    \mathbb{Z}_2^2 = \left\{00, 01, 10, 11\right\}
\end{equation*}
The \hl{operation is component-wise addition modulo 2, which is equivalent to the XOR operation.} So throughout Simon, whenever we've written $x \oplus a$, we've actually been performing addition in the group $\mathbb{Z}_2^n$. That's why Simon's promise:
\begin{equation*}
    f(x) = f(x \oplus a)
\end{equation*}
Can be interpreted as saying: \hl{$f$ possesses a hidden periodic symmetry under addition by $a$.}

\highspace
\begin{flushleft}
    \textcolor{Green3}{\faIcon{question-circle} \textbf{Where does Fourier analysis enter in?}}
\end{flushleft}
That's the fun part: we've already been doing Fourier analysis! Recall the formula we've used repeatedly:
\begin{equation*}
    H^{\otimes n} \ket{x} = \frac{1}{\sqrt{2^n}} \sum_{y \in \mathbb{Z}_2^n} (-1)^{x \cdot y} \ket{y}
\end{equation*}
The factors $(-1)^{x \cdot y}$ are the \textbf{Fourier phases} appropriate to the group $\mathbb{Z}_2^n$. So the final Hadamards \textbf{aren't just an arbitrary trick} that happens to produce interference. Mathematically, $H^{\otimes n}$ is \textbf{performing the Fourier transform} for the particular algebraic structure in which Simon's period lives.

\highspace
Thus, for bitstrings with XOR, this is called the \definition{Fourier transform over $\mathbb{Z}_2^n$}, or equivalently the \definition{Walsh-Hadamard transform}.

\highspace
\begin{flushleft}
    \textcolor{Green3}{\faIcon{book} \textbf{Simon as a hidden subgroup problem}}
\end{flushleft}
There's an even more abstract interpretation, but we only need the basic idea. Consider the set:
\begin{equation*}
    H = \left\{
        0^{n}, a
    \right\}
\end{equation*}
It forms a \textbf{small subgroup} of $\mathbb{Z}_2^n$, and Simon's function is constant on pairs of the form:
\begin{equation*}
    \left\{
        x, x \oplus a
    \right\}
\end{equation*}
These pairs are precisely the \textbf{cosets} (i.e., shifted copies) associated with that hidden subgroup. So, \hl{Simon hides the subgroup $H$ inside the equality structure of $f$.} The quantum algorithm uses Fourier sampling to obtain information that lets us reconstruct that hidden structure.

\highspace
In other words, the subgroup is the set of XOR shifts that don't change the function's output $f$. Obviously the trivial subgroup $\left\{0^n\right\}$ is always hidden, but the promise of Simon's problem is that there is a \textbf{non-trivial} hidden subgroup $\left\{0^n, a\right\}$, and the quantum algorithm is able to find it efficiently.

\begin{examplebox}[: Simon's hidden subgroup]
    Suppose Simon's secret is $a = 01$. The possible inputs are:
    \begin{equation*}
        \left\{00, 01, 10, 11\right\}
    \end{equation*}
    Simon promises that:
    \begin{equation*}
        f(x) = f(x \oplus a)
    \end{equation*}
    So the function is constant on the pairs:
    \begin{equation*}
        \underbrace{\left\{00, 01\right\}}_{
            00 \oplus 01 = 01
        },\quad \underbrace{\left\{10, 11\right\}}_{
            10 \oplus 01 = 11
        }
    \end{equation*}
    Therefore, the domain is partitioned.

    \highspace
    Now, the \textbf{hidden subgroup} is:
    \begin{equation*}
        H = \left\{00, 01\right\}
    \end{equation*}
    And it is special because if we combine its elements using XOR, we never leave the original set:
    \begin{align*}
        00 \oplus 00 &= 00,\\
        00 \oplus 01 &= 01,\\
        01 \oplus 00 &= 01,\\
        01 \oplus 01 &= 00
    \end{align*}
    So $H$ has a closed little structure under XOR. The other coset is obtained by taking an element outside $H$, for example $10$, and XORing it with every element of $H$:
    \begin{equation*}
        10 \oplus H
        =
        \left\{
            10 \oplus 00,\,
            10 \oplus 01
        \right\}
        =
        \left\{10,11\right\}
    \end{equation*}
    Therefore, the two cosets are:
    \begin{equation*}
        \underbrace{\left\{00,01\right\}}_{H}
        \qquad\text{and}\qquad
        \underbrace{\left\{10,11\right\}}_{10\oplus H}
    \end{equation*}
    The \textbf{second coset is not a subgroup}. In particular, it does not contain the identity element $00$ and it is not closed under XOR. For example:
    \begin{equation*}
        10 \oplus 11 = 01
        \notin \left\{10,11\right\}
    \end{equation*}
    \textcolor{Green3}{\faIcon{balance-scale} \textbf{Subgroup vs. Coset.}} A \textbf{subgroup} is a \hl{group inside the larger group}, while a \textbf{coset} is a \hl{shifted copy of a subgroup.} In this case, $H$ is a subgroup because it contains the identity $00$, and XORing any two elements of $H$ keeps us inside $H$. In contrast, $10 \oplus H$ XORs with \textbf{every element of $H$} and produces:
    \begin{equation*}
        10 \oplus H = \left\{10,11\right\}
    \end{equation*}
    It is not a subgroup because it does not contain the identity $00$, but it is a coset because it is a shifted copy of $H$.
\end{examplebox}

\highspace
\textcolor{Green3}{\faIcon{question-circle} \textbf{Why is it called hidden?}} Simple, because we \textbf{don't know} $a$. If we knew $a$, we could easily find the subgroup $H = \left\{0^n, a\right\}$ instantly. But the oracle doesn't tell us $H$. Instead, we can only query $f$, and the structure of $f$ indirectly reveals that:
\begin{equation*}
    f(x) = f(x \oplus a)
\end{equation*}
Then:
\begin{equation*}
    H = \left\{0^n, a\right\}
\end{equation*}
Because \textbf{finding this subgroup is essentially equivalent to finding $a$}.

\highspace
\begin{flushleft}
    \textcolor{Red2}{\faIcon{exclamation-triangle} \textbf{Ordinary Period Finding}}
\end{flushleft}
Suppose instead that our inputs aren't $n$-bit strings under XOR. Imagine:
\begin{equation*}
    f: \mathbb{Z}_{N} \to \mathbb{Z}_{N}
\end{equation*}
With an ordinary modular period $r$:
\begin{equation*}
    f(x) = f((x + r) \mod N)
\end{equation*}
For example, suppose $N=12$ and $r=3$. Then we'd have patters like:
\begin{gather*}
    f(0) = f(3) = f(6) = f(9),\\
    f(1) = f(4) = f(7) = f(10),\\
    f(2) = f(5) = f(8) = f(11)
\end{gather*}
Compare that with Simon:
\begin{equation*}
    f(x) = f(x \oplus a)
\end{equation*}
The conceptual problem is extremely similar: \textbf{there is a hidden period, we need to recover it}. But the underlying mathematical structure has changed. Simon has $\mathbb{Z}_{2}^n$ as its domain, while the ordinary period finding problem has $\mathbb{Z}_{N}$ as its domain. And \hl{this is why we need something more general than Hadamards}.

\highspace
\textcolor{Green3}{\faIcon{question-circle} \textbf{Why isn't $H^{\otimes n}$ enough anymore?}} Because the Hadamard only introduces phases of $+1$ and $-1$. We can see that directly from $(-1)^{x \cdot y}$. Those phases are perfect for binary/XOR structure. But suppose the period lives in something like $\mathbb{Z}_{N}$, where $N$ is not a power of 2. Then we need richer phases corresponding to the $N$-th roots of unity:
\begin{equation*}
    1, \quad e^{2\pi i / N}, \quad e^{4\pi i / N}, \quad \ldots
\end{equation*}
That motivates the general \textbf{Quantum Fourier Transform}:
\begin{equation*}
    QFT_{N} \ket{x} = \frac{1}{\sqrt{N}} \sum_{y=0}^{N-1} e^{2\pi i xy / N} \ket{y}
\end{equation*}
We'll explore this in more detail in the next sections, but for now, compare it with Simon:
\begin{equation*}
    H^{\otimes n} \ket{x} = \frac{1}{\sqrt{2^n}} \sum_{y \in \mathbb{Z}_2^n} (-1)^{x \cdot y} \ket{y}
\end{equation*}
They have the \textbf{same broad structure}: a \hl{superposition over $y$ with phases depending on $x$ and $y$.} The difference is that Simon's underlying group is $\mathbb Z_2^n$, for which the Fourier transform happens to be simply $H^{\otimes n}$.